\documentclass[11pt]{article}
\usepackage[margin=1in]{geometry}
\usepackage[T1]{fontenc}
\usepackage{amsmath,amssymb,mathtools}
\usepackage{booktabs,array,enumitem,caption}
\usepackage{xcolor,hyperref}
\sloppy
\hypersetup{colorlinks=true, linkcolor=blue!50!black, urlcolor=blue!50!black}
\newcommand{\code}[1]{\texttt{#1}}
\newcommand{\file}[1]{\texttt{#1}}
\setlist{nosep, leftmargin=1.4em}

\title{AutoBayes brms-parity loop: Tier F review}
\author{Status at the Tier F boundary, prepared for Andrew}
\date{25 October 2026}

\begin{document}
\maketitle

\noindent This is the Tier F boundary review. Since 21 October you asked for six things: continue with Tier F, benchmark
against brms, keep a log of brms's bugs for a later pull request, check that the library reads brms's syntax without
extraneous parts, write unit tests for all code before it merges, and move the agent scaffold out of the repository. All of
Tier F is on \code{main}, and every one of the 45 response families brms 2.23 documents now has a verified model.
Section~\ref{sec:brms} is what brms gets wrong, Section~\ref{sec:syntax} answers the syntax question, and
Section~\ref{sec:decisions} is what needs you.

\section{What was built}

\begin{table}[tbp]
\small
\begin{center}\begin{tabular}{@{}>{\raggedright\arraybackslash}p{0.17\textwidth}>{\raggedright\arraybackslash}p{0.07\textwidth}>{\raggedright\arraybackslash}p{0.52\textwidth}>{\raggedright\arraybackslash}p{0.13\textwidth}@{}}
\toprule
Feature & PR & What a user can now write & Diverg. \\
\midrule
F01a multinomial & \#92 & \code{ab.multinomial("cbind(m1, m2, m3) | trials(size) \textasciitilde{} x + (1 | g)", ...)}, and \code{ab.dirichlet\_multinomial} & D136--D137 \\
F01b simplex & \#95 & \code{ab.dirichlet("cbind(s1, s2, s3) \textasciitilde{} x + (1 | g)", ...)}; \code{ab.logistic\_normal(..., refcat="s2")} & D138--D139 \\
F02a CAR & \#93 & \code{car(W, gr = site, type = "bym2")} (also escar, esicar, icar), \code{data2=\{"W": W\}} & D140--D141 \\
F02b SAR & \#87 & \code{sar(M, type = "lag")} and \code{"error"} & D142--D143 \\
F03a fcor & \#91 & \code{fcor(V)}, a fixed residual covariance & D144--D145 \\
F03b unstr, cosy & \#96 & \code{unstr(time = tt, gr = subj)}, \code{cosy(time = tt, gr = subj)} & D146--D147 \\
F04 wiener & \#88, \#90 & \code{ab.wiener("rt | dec(dec) \textasciitilde{} x + (1 | g)", ...)}, with formulas on \code{bs} & D148--D151 \\
F05 cox & \#101, \#103 & \code{ab.cox("time | cens(c) + bhaz(df = 7) \textasciitilde{} x + (1 | g)", ...)} & D152--D155 \\
F06a se, rate & \#98 & \code{y | se(sdx) \textasciitilde{} 1 + (1 | g)}, \code{se(sdx, sigma = TRUE)}; \code{cnt | rate(exposure) \textasciitilde{} x} & D156--D157 \\
F06b thres & \#97 & \code{ab.cumulative("ord | thres(gr = z) \textasciitilde{} x + (1 | g)", ...)}, \code{thres(3, gr = z)} & D158--D159 \\
F06c subset & \#99 & \code{["y | subset(sub1) \textasciitilde{} x + (1 | g)", "y2 | subset(sub2) \textasciitilde{} x"]} & D160--D161 \\
F07 summaries & \#102 & \code{ab.fixef}, \code{ranef}, \code{coef}, \code{VarCorr}, \code{summary}, \code{prior\_summary}, \code{residuals}, \code{predictive\_error} & D162--D165 \\
F08 syntax fixes & \#105 & ordered factors (\code{orderedz.L}, \code{.Q}); \code{(x | g)}; \code{gr(g, cor = FALSE)}; \code{trials(15)} & D168--D169 \\
F09 names, messages & \#106 & \code{from autobayes.session import fit\_record}: twelve internal names leave the top level; messages without milestone names & D172--D175 \\
\bottomrule
\end{tabular}\end{center}
\caption{\small The Tier F features, verified on 45 frozen brms cases: \#83, \#89 and \#101 froze 38 before the features
began, and the features added 7.}\label{tab:built}
\end{table}

\noindent Table~\ref{tab:built} lists the features. Seven more pull requests changed how the loop works rather than what the
library does:
\begin{itemize}
\item \textbf{Diff coverage (\#86).} A feature merges only if every library line it changes runs under the tests. This is
  how I read your ``unit tests for all code before merging''. The rule is: 100\% of the changed executable lines under
  \file{autobayes/}, with exemptions only as a pragma that states its reason. \#94 then covered the 13 lines \#77 had left
  unrun.
\item \textbf{The syntax corpus (\#85).} 248 brms formulas, each read by brms and by AutoBayes; a test fails if any reading
  changes (Section~\ref{sec:syntax}).
\item \textbf{The upstream log (\#81).} Every problem in brms, Stan Math, loo and posterior gets a standalone R
  reproduction. It is kept out of git since \#82, as you asked.
\item \textbf{The benchmark (\#84, \#104)}: Section~\ref{sec:speed}.
\item \textbf{The tidy (\#100).} The agent scaffold moved to \href{https://github.com/andrew-dou/autobayes-scaffold}{andrew-dou/autobayes-scaffold}
  (private), and the milestone history to \file{docs/history/}. The top level is now \file{autobayes}, \file{docs},
  \file{examples}, \file{oracle}, \file{tests}, \file{tools} and the packaging files.
\end{itemize}

\paragraph{How close.} I measured with the gates' own comparisons on \code{main} after F08 (\code{203f5db}). Every gate
passed on all 45 Tier F cases, with two exceptions described in Section~\ref{sec:checks}: the two bym2 G1 failures,
which are expected, and wiener's expected values, which are accepted against references. The worst relative errors over
the Tier F cases:
\begin{itemize}
\item $6.7\times10^{-16}$ in G1;
\item $1.9\times10^{-9}$ in G2's log density (\code{cox\_uncensored\_default\_priors}), a fifth of the tolerance. I have
  not traced its cause.
\item $1.2\times10^{-8}$ in G2's gradient (\code{wiener\_distributional\_bs}). This is brms's error: Stan's
  \code{wiener\_lpdf} truncates its series, and AutoBayes's density is exact to $4\times10^{-15}$ against a 60-digit
  evaluation (D150).
\item $6.1\times10^{-10}$ in G3 (\code{poisson\_car\_bym2}), which comes through brms's inexact bym2 scaling factor;
\item $8.8\times10^{-16}$ in predictions;
\item $2.2\times10^{-15}$ between the JAX and PyTensor densities;
\item $1.4\times10^{-14}$ in the post-processing tables (\code{gauss\_sar\_lag\_post}).
\end{itemize}
Both G2 figures are also the worst over all cases, Tier A to F, outside the cases whose gates are expected to fail.
Everything is inside the unchanged tolerances.
The full non-sampling suite on \code{main} after F09 (\code{5cfcb0b}) passed: 6798 tests passed and none failed, with 193 skipped and 28 expected failures.
\code{oracle/regen.py check --all} regenerated all 216 frozen fixtures from R on \code{a9b5cf3}, and every one was
identical. F08's new case, regenerated on its own afterwards, was identical too.

\section{What brms gets wrong}
\label{sec:brms}

The upstream log now has 52 entries: 41 in brms (one of them not a bug after all), 6 in Stan Math, 3 in loo, and one each
in posterior and R. Tier F added 21 of them (U32--U52), each with a script that prints the library's value next to the
exact one and exits non-zero while the problem is present. The ones that give a wrong value, or stop where a value
exists:
\begin{itemize}
\item \textbf{SAR models (U32, U33, U36).} A row-standardised weight matrix, the usual kind, cannot be built: R's
  \code{eigen()} returns complex eigenvalues and rstan refuses them. \code{posterior\_epred()} of an error-type SAR model
  stops. Both SAR log densities leave out the Gaussian constant $-\tfrac{N}{2}\log 2\pi$.
\item \textbf{The Wiener mean (U34).} \code{posterior\_epred()} uses the relative starting point where the absolute one
  belongs, so a mean response time can come out negative. AutoBayes's mean is accepted against a 50-digit reference
  (\#90).
\item \textbf{\code{bayes\_R2()} with \code{unstr()} or ARMA terms (U41).} When the data are not sorted by group and time,
  brms pairs the responses with the means in a different order. On one frozen case its value is up to 23.5\% off.
\item \textbf{Cox models (U45--U47).}
  \begin{itemize}
  \item Interval censoring gives a log density of $-\infty$ at every point.
  \item \code{trunc(lb = L)} normalises each row by its own survival $S(y)$ instead of $S(L)$, which cancels the row's
    survival term.
  \item The default Intercept prior is centred at $\log(\text{median time})$, so it moves with the unit of time although
    the intercept does not.
  \item \code{posterior\_linpred()} on new rows stops without the time column, which it does not use (U50).
  \end{itemize}
  AutoBayes refuses the first two, records the third, and predicts without the time column.
\item \textbf{\code{factor()} and \code{ordered()} written in a formula (U51).} New rows that lack one of the levels fail
  with ``non-conformable arguments'', and so does a single new row. F08 found it while freezing its ordered-factor case.
  \code{lm()} handles the same formula through its recorded levels. AutoBayes keeps the levels it was fitted with (for
  declared categoricals and for \code{ordered()}), so new rows need not hold every level.
\end{itemize}
The rest are inputs brms should refuse and does not:
\begin{itemize}
\item disconnected graphs in \code{car()}, and a diagonal entry counted as an edge (U38, U39);
\item SAR bounds of $(\lambda_{\min}, \lambda_{\max})$ where $(1/\lambda_{\min}, 1/\lambda_{\max})$ belongs (U35);
\item \code{dec()} reading any non-zero number as the upper boundary (U40);
\item negative multinomial counts (U42), and per-category formulas silently dropping the response formula's terms (U43);
\item a zero standard error in \code{se()} (U44);
\item \code{bhaz(df = 7, knots = c(1.5, 3))}, which silently ignores \code{df} (U49).
\end{itemize}
Three are inexact:
\begin{itemize}
\item bym2's scaling factor is INLA's perturbed value, not Riebler's exact one: $1.5\times10^{-8}$ to $2.3\times10^{-7}$
  low (U37);
\item posterior's \code{rhat()} and ESS return numbers for chains of two or three iterations (U48);
\item R's \code{contr.poly()}, which brms uses for ordered factors, drifts from the polynomial contrasts as the levels
  grow: $1.6\times10^{-12}$ at 17 levels, and from 23 levels on its top columns are not polynomial at all (U52). R core
  answered a related report in 2002 that such polynomials cannot be represented accurately, but the exact ones can be:
  the script computes them to $10^{-15}$. AutoBayes builds them from the definition and refuses more than 16 levels
  (D168).
\end{itemize}
Your rule of 24 October settled how AutoBayes treats these: follow brms where it is logical, and otherwise do the sensible
thing and record the difference. One case went the other way. For bym2, AutoBayes keeps INLA's definition, computed
exactly, because the exact factor would buy nothing practical and would make the model unverifiable against brms (D140).

The drafts in \file{upstream-drafts/} are ready to post: the maintainer note, the pull-request plan and the log's PDF
(Section~\ref{sec:decisions}).

\section{Speed}
\label{sec:speed}

The benchmark of 25 October (\#104) fitted 18 models with both libraries: the same formulas, priors and data, 4 chains of
1000 warm-up and 1000 kept draws, and target acceptance 0.9.
\begin{itemize}
\item \textbf{End to end}, compiling and sampling, brms took 0.8 to 18.1 times as long (median 6.5).
\item \textbf{Its sampler alone} took 0.6 to 13.2 times as long (median 3.8).
\item \textbf{Where the gain comes from.} The two samplers take about the same number of leapfrog steps. The gain is the
  gradient: on the median model, B11, one gradient costs 32.8\,$\mu$s in JAX and 206.9\,$\mu$s in Stan.
\item \textbf{Where brms wins.} brms is faster on the horseshoe model (B8) and the exact Gaussian process (B12).
\item \textbf{Agreement.} The largest difference of posterior means over their combined Monte Carlo error was
  $|z| = 2.71$ across all 18 models.
\end{itemize}

\section{Does AutoBayes read brms's syntax?}
\label{sec:syntax}

You asked on 24 October whether the code really matches brms's syntax, without extraneous parts. The audit (\#85) read 248
formulas with both libraries and compared what each builds, \code{make\_standata()} against AutoBayes's design. The
corpus is now a test that fails if any reading changes.
\begin{itemize}
\item \textbf{Before F08:} 114 of the 116 formulas AutoBayes accepted built brms's data. One differed by design (D9), and
  one was a silent misreading: an ordered factor was treatment-coded where brms uses polynomial contrasts. Of the 132
  refused, 16 were refused without naming the construct.
\item \textbf{After F08:} 134 formulas match, 113 are refused with the construct named, and the one documented difference
  remains. Nothing is misread. Ordered factors now use R's polynomial contrasts, built from their definition and checked
  against brms on a new frozen case. \code{(x | g)} now means \code{(1 + x | g)}, as in brms.
\item \textbf{Extraneous parts (F09).} Twelve names of the scaffold's session records and of the term model's internals (\code{fit\_record},
  \code{environment}, \code{ModelSpec}, \code{parse\_formula} and others) left \code{\_\_all\_\_}. Each is still
  importable from its module, and \code{ab.<name>} still works with a deprecation warning. \code{describe()} names the
  family instead of M0, M1, M2, and about 145 refusals say ``not supported'' instead of ``in this slice'' and cite current
  docs instead of milestone decisions. No refusal changed what it refuses, and the corpus totals are unchanged. The
  call-surface differences that had no divergence row now have one: declared factors (D172), how priors cover a class
  (D173), rows with missing values (D174), and \code{fit()}'s defaults and the posterior's names (D175).
  Section~\ref{sec:decisions} lists the four changes F09 proposes but did not make.
\end{itemize}
What remains refused is listed construct by construct in \file{docs/current-capabilities.md}. The largest gaps are
\code{s()} and \code{t2()} smooths, three-way interactions, \code{poly()}, and many family--context pairs: a family
added late often takes the common terms but not \code{cens()}, \code{weights()} or a distributional formula.

\section{What combining the features found}

\begin{itemize}
\item \textbf{One \code{data2}.} CAR, SAR and fcor each read a matrix from \code{data2}, and fcor's and SAR's branches
  conflicted in 47 places. Their agents unified them into one mechanism used by all three. Since then each agent rebases
  onto \code{main} itself and reruns its own checks before reporting, which ended the large conflicts at merge time.
\item \textbf{Lost lines in rebases.} Resolving conflicts silently dropped three lines: a \code{continue}, which let
  \code{sar()} fall through to the autocorrelation parser, a \code{@property}, and a function's \code{return}. The reruns
  after each rebase caught all three before merge. Agents now check every function both sides touched against
  \code{main}.
\item \textbf{The corpus.} Re-recording all 248 entries rewrote unrelated ones whenever brms's message for some other
  construct changed. A helper now re-records only the entries a feature changes.
\item \textbf{A hash-seed-dependent \code{describe()}.} The order of two lines of one hurdle-cumulative model changes with
  \code{PYTHONHASHSEED}. Several agents found it. F09 traced it to Bambi, not brms: \code{Family.auxiliary\_parameters} is a Python set, and
  \code{bmb.Model} creates the parameters' variables in its iteration order. It reached past \code{describe()}: in 25
  cases (learned Student-t \code{nu}, the hurdle and zero-inflated families, wiener and others) the free variables came in
  a different order in each process. An existing reordering hid it everywhere but that one \code{describe()}. A seeded
  fit of those models was then probably not reproducible across processes; this is inferred, since nothing was sampled.
  The family handed to Bambi now lists them in its likelihood's order, tested in fresh interpreters under two hash
  seeds.
\end{itemize}

\section{How the claims were checked, and their limits}
\label{sec:checks}

\begin{itemize}
\item \textbf{Per feature.} The gate tool ran G1--G6 on the tree that merged. Diff coverage reached 100\% on every feature
  from \#87 on, and the corpus test passed. F01a's gate record was lost from the ledger in its rebase; its gate log
  (GATES GREEN on \code{f138205}) survives, and the record is restored. Each agent ran the
  full non-sampling suite on its own branch. I re-ran the focused tests after any rebase.
\item \textbf{Where brms's number is not the target.} Four Tier F cases are not checked against brms alone:
  \begin{itemize}
  \item the two bym2 cases, frozen with G1 expected to fail: brms's scaling factor is $1.7\times10^{-10}$ from its own
    definition, beyond G1's $10^{-12}$; every other gate passes (D140);
  \item the two wiener cases, whose expected values are accepted against 50-digit references instead of brms's wrong
    means (D148).
  \end{itemize}
\item \textbf{Exceptions.} The full suite ran on \code{main} only at boundaries, not after every merge. Between
  boundaries a merge rested on the agent's own full suite and my focused re-run.
\item \textbf{Limits.} Nothing in the loop samples. The only fits are the benchmark's 18 models. No Tier F family has been
  fitted outside its frozen points, so whether NUTS samples a CAR, SAR or cox posterior well is untested.
\end{itemize}

\section{Process and cost}

Agent tokens reported, in millions:
\begin{itemize}
\item fixtures: 1.20 (three agents);
\item features: F01a 0.94, F01b 0.92, F03b 0.90, F02a 0.84, F04 0.81, F08 0.73, F05 0.72, F06a 0.72, F07 0.67,
  F06b 0.66, F06c 0.65, F03a 0.63, F02b 0.63, F09 0.53;
\item the loop's own tools: upstream log 0.97, syntax audit 0.57, stage benchmark 0.43, tidy 0.27,
  diff coverage 0.25, \#77's lines 0.23;
\item about 14.3M in all, against Tier E's 4.6 for six features.
\end{itemize}
Up to ten agents ran at once. Two things limited the pace: two full suites at a time on this 16\,GB machine, and, on
24 October, a disk at 8.9\,GB free with swap almost full. Network outages, an expired login and stalled streams stopped
agents several times; each resumed from its own context without losing work. I paused three agents once, at 56\% of a five-hour window, and
resumed them after it reset.

\section{Decisions}
\label{sec:decisions}

\begin{enumerate}
\item \textbf{Posting to brms.} The maintainer note, the pull-request plan and the log's PDF are in
  \file{upstream-drafts/}. They go out under your account. I suggest posting the note alone first, as an issue, and
  writing code only after Paul says how he wants it split.
\item \textbf{The scaffold repository} is private. Make it public? Either way, AutoBayes's public history before \#100
  still contains the scaffold, and \#81's history contains the upstream log.
\item \textbf{Two withheld results, under your 24 October rule.} Both are places where brms is wrong and AutoBayes computes
  the right answer, which your rule now covers:
  \begin{itemize}
  \item the truncated count mean $E[Y \mid Y \ge lb]$, where brms gives $E[Y \mid Y \ge lb + 1]$ (D121);
  \item stacking weights at the exact optimum, where loo stops $1.0\times10^{-8}$ short (D126).
  \end{itemize}
  Unless you object, I will release both, each with its divergence row and a 50-digit reference.
\item \textbf{Splines.} \code{s()} and \code{t2()} are the largest syntax gap left. B04's phase 1 (a basis port checked
  against R) is on its branch, waiting on the licensing question: GPL'd mgcv, and whether a port written from mgcv's
  documentation and Wood's papers is acceptable.
\item \textbf{CI billing} is still blocking hosted runs; every merge has been verified locally.
\item \textbf{Four call-surface changes F09 proposes.} Each departs from today's refusal toward brms; none is made yet.
  \begin{itemize}
  \item Accept an undeclared pandas \code{Categorical}, its first category the reference, as R does for a factor; and
    accept a declared boolean column, which is refused today although brms codes it \code{FALSE}/\code{TRUE} (D172).
  \item Allow a prior on a whole class plus overrides for single coefficients, brms's \code{coef} rows, in place of
    today's mapping that must name every term (D173).
  \item An opt-in to drop rows with missing values and report them, as zero-trial rows already are (D174).
  \item Change \code{target\_accept}'s default from 0.9 to brms's 0.8 (D175).
  \end{itemize}
  I suggest the first three, under your rule, and keeping 0.9: it costs a little speed and gives fewer divergent
  transitions, and anyone matching brms can pass 0.8.
\item \textbf{Tier G.} Candidates, in the order I would take them:
  \begin{enumerate}
  \item[(a)] the common syntax still refused: three-way interactions, interactions with transforms, and \code{poly()},
    which is built like the ordered-factor contrasts;
  \item[(b)] the family--context matrix: \code{cens()}, \code{weights()}, distributional formulas and group-level terms
    for the families that do not yet take them;
  \item[(c)] the two releases in item 3, and whichever changes in item 6 you accept;
  \item[(d)] a first real fit of each Tier F family, on the GPU or the cluster, which needs your approval for compute;
  \item[(e)] splines, once item 4 is decided.
  \end{enumerate}
  At Tier F's rate, about 0.75M tokens a feature, (a) and (b) together are perhaps 8--12M.
\end{enumerate}

\end{document}
